\documentclass[12pt]{article}
\usepackage{tikz}
\usetikzlibrary{shapes}

\usepackage[utf8]{inputenc}
\usepackage{amsmath,amsfonts,amssymb,amsthm,mathrsfs,graphicx,
	setspace,lscape,longtable,multirow,chngcntr,
	color,bm,booktabs,float,enumitem,
	titlesec,apptools,makecell,array}
\usepackage[T1]{fontenc}
\usepackage[colorlinks=true,linkcolor=blue,urlcolor=blue,citecolor=blue]{hyperref}
\usepackage{xurl} % Allows URLs to break at any alphanumeric character
\usepackage[labelfont=bf]{caption,subcaption}
\usepackage[table, dvipsnames]{xcolor}
\usepackage[bottom]{footmisc}
\usepackage[page,title]{appendix}
\usepackage{natbib}
\usepackage[top=1.25in, bottom=1.25in, left=1.25in, right=1.25in]{geometry}
\usepackage{ragged2e}
\usepackage{placeins}

\renewcommand{\today}{\ifcase \month \or January\or February\or March\or %
	April\or May \or June\or July\or August\or September\or October\or November\or %
	December\fi \ \number \year}

%%%%%%%%%%%%%%%% Define math environments %%%%%%%%%%
\AtAppendix{\counterwithin{theorem}{section}}
\newtheorem{theorem}{Theorem}
\newtheorem{algorithm}{Algorithm}
\newtheorem{axiom}{Axiom}
\newtheorem{case}[theorem]{Case}
\newtheorem{claim}{Claim}
\newtheorem{conclusion}{Conclusion}
\newtheorem{condition}[theorem]{Condition}
\newtheorem{conjecture}{Conjecture}
\AtAppendix{\counterwithin{corollary}{section}}
\newtheorem{corollary}{Corollary}
\AtAppendix{\counterwithin{definition}{section}}
\newtheorem{definition}{Definition}
\AtAppendix{\counterwithin{lemma}{section}}
\newtheorem{lemma}{Lemma}
\AtAppendix{\counterwithin{assumption}{section}}
\newtheorem{assumption}{Assumption}
\AtAppendix{\counterwithin{proposition}{section}}
\newtheorem{proposition}{Proposition}
\AtAppendix{\counterwithin{remark}{section}}
\newtheorem{remark}{Remark}
\newtheorem{exmp}{Example}

%%%%%%%%%%%%%%%%% Common math %%%%%%%%%%%%%%%%%%%%%%%%%
\newcommand{\iid}{\stackrel{\mathrm{iid}}{\sim}}
\newcommand{\prob}{\mathbb{P}}
\newcommand{\expec}{\mathbb E}
\newcommand{\var}{\mathrm{var}}
\newcommand{\cov}{\mathrm{cov}}
\newcommand{\corr}{\mathrm{corr}}
\newcommand{\diff}{\mathrm d}
\newcommand{\intdiff}{\,\mathrm d}
\DeclareMathOperator*{\subjecto}{ subject\ to\quad }
\DeclareMathOperator*{\maximize}{\mathrm maximize\quad }
\DeclareMathOperator*{\minimize}{\mathrm minimize\quad }
\DeclareMathOperator*{\argmin}{\displaystyle arg\,min}
\DeclareMathOperator*{\argmax}{\displaystyle  arg\,max}
\DeclareMathOperator*{\normal}{\mathcal{N}}
\DeclareMathOperator{\Var}{Var}
\DeclareMathOperator{\Cov}{Cov}
\newcommand{\Phinv}{\Phi^{-1}}
\newcommand{\E}{\mathbb{E}}

% Needed for "code like" text
\def\code#1{\texttt{#1}}

%%%%%%%%%%%%%%%%%%% Paper-specific commands %%%%%%%%%%%
%% Add your custom \newcommand definitions here.
%% Examples:
%%   \newcommand\bmR{\bm R}        % bold return vector
%%   \newcommand\SR{\mathrm{SR}}   % Sharpe ratio
%%   \newcommand\BF{\mathrm{BF}}   % Bayes factor
%%%%%%%%%%%%%%%%%%%%%%%%%%%%%%%%%%%%%%%%%%%%%%%%%%%

%%%% Hide Columns %%%%%%%%%%%
\newcolumntype{H}{>{\setbox0=\hbox\bgroup}c<{\egroup}@{}}
%%%% Math Columns %%%%%%%%%%%
\newcolumntype{C}{>{$}c<{$}}
\newcolumntype{R}{>{$}r<{$}}
\newcolumntype{L}{>{$}l<{$}}

\renewcommand{\thesubfigure}{\Alph{subfigure}:}
\captionsetup[subfloat]{font=footnotesize,labelformat=simple,labelsep=space}

\singlespacing
\title{
	{\textbf{[REMOVE] Your Paper Title Here}}\thanks{{\scriptsize
	[REMOVE] We gratefully acknowledge comments and suggestions from \ldots
}}
}

\author{
{\normalsize [REMOVE] First Author\thanks{{\scriptsize
[REMOVE] Department of Finance, University Name; \url{first.author@university.edu}}}}
\and
{\normalsize [REMOVE] Second Author\thanks{{\scriptsize
[REMOVE] Department of Economics, University Name; \url{second.author@university.edu}}}}
}

\date{{\small   \today }}

\evensidemargin -0.1in \oddsidemargin -0.1in \textwidth 6.7in
\topmargin -0.5in \textheight 9in
\setlength{\footnotesep}{10.5pt} \linespread{2}
\renewcommand{\baselinestretch}{1.3}
\setlength{\footskip}{40pt}

\linespread{1.5}


\begin{document}
\maketitle

\vspace{-.75cm}

\begin{abstract}
\singlespacing
\noindent
{[REMOVE] This paper examines the post-earnings announcement drift across 40 international equity markets. Combining firm-level earnings surprises with daily return data from 1990 to 2023, we document substantial cross-country variation in the magnitude and persistence of the drift. In developed markets, the average cumulative abnormal return over the 60-day post-announcement window is 1.8\%, while emerging markets exhibit a drift of 3.4\%. Transaction costs absorb between one-third and two-thirds of the gross premium depending on the market.}
\end{abstract}
\bigskip
\singlespacing{
\noindent \emph{Keywords:} [REMOVE] Keyword One; Keyword Two; Keyword Three.}
\medskip

\noindent \emph{JEL Classification Codes:} [REMOVE] G12; C12.
\bigskip
\thispagestyle{empty}

\pagebreak
\setcounter{page}{1}
\onehalfspacing
\setstretch{1.4}
\setlength{\belowdisplayskip}{5pt}
\setlength{\belowdisplayshortskip}{5pt}
\setlength{\abovedisplayskip}{5pt}
\setlength{\abovedisplayshortskip}{5pt}

%%%%%%%%%%%%%%%%%%%%%%%%%%%%
%%% Sections in the main body
%%%%%%%%%%%%%%%%%%%%%%%%%%%%

%%%%%%%%%%%%%%%%%%%%%%%%%%%%
% INTRODUCTION SECTION
%%%%%%%%%%%%%%%%%%%%%%%%%%%%
\clearpage

%% BEGIN:introduction
\section{Introduction}\label{sec:intro}

[REMOVE] The post-earnings announcement drift (PEAD) is one of the most robust anomalies in equity markets, first documented by \citet{FamaFrench_1992} and studied in the context of multifactor models by \citet{FamaFrench_1993}. [REMOVE] Despite four decades of research, the source and persistence of the drift remain contested: explanations range from investor underreaction to earnings news to risk-based channels tied to information uncertainty. [REMOVE] A growing body of international evidence suggests that institutional features such as disclosure frequency, analyst coverage, and short-selling constraints shape both the magnitude and the speed of price adjustment.

[REMOVE] We contribute to this literature in three ways. [REMOVE] First, we construct a standardized earnings surprise measure across 40 countries using a common methodology, eliminating cross-study differences in how surprises are defined. [REMOVE] Second, we decompose the drift into anticipation, announcement, and post-announcement windows to isolate where the return accrues. [REMOVE] Third, we evaluate the net-of-cost profitability of drift-based strategies under realistic transaction cost assumptions, building on the framework in \citet{FamaFrench_2015}.

[REMOVE] The remainder of this paper proceeds as follows. Section~\ref{sec:data} describes the data sources, sample construction, and variable definitions. Section~\ref{sec:methodology} develops the econometric framework for measuring abnormal returns around earnings announcements. Section~\ref{sec:results} presents the main empirical results. Section~\ref{sec:conclusion} concludes.

%% END:introduction


%%%%%%%%%%%%%%%%%%%%%%%%%%%%
% DATA SECTION
%%%%%%%%%%%%%%%%%%%%%%%%%%%%

%% BEGIN:data
\section{Data}\label{sec:data}

[REMOVE] This section describes the data sources, sample construction, and variable definitions used throughout the paper. [REMOVE] Our sample covers 40 equity markets from January 1990 to December 2023 and includes all common stocks with at least 120 trading days per year and a market capitalization above the 20th percentile of their domestic exchange. [REMOVE] Table~\ref{tab:example_descriptive} reports summary statistics for the main variables.

%%% --- TABLE 1: Panel Results with T-Stats (Archetype 1) --- %%%
%%% The workhorse results table: multi-column comparison with panels and t-stats.
%%% Pattern: \scalebox{0.75}, booktabs, \cmidrule(lr), t-stats on sub-rows, \addlinespace.
\begin{table}[!ht]
\caption{[REMOVE] Post-Announcement Drift by Region: Value-Weighted Portfolios.}
\begin{spacing}{1}
{\footnotesize
[REMOVE] $\mu$ is mean monthly return (\%), $\alpha$ is Fama-French three-factor alpha (\%). Columns compare developed markets (DM) and emerging markets (EM). $t$-statistics (Newey-West, lags $= \lfloor T^{0.25} \rfloor$) in parentheses. Sample: 1990-01 to 2023-12, $T$=408.}
\end{spacing}
\vspace{-4mm}
\begin{center}
\label{tab:example_panel}
\vspace{2mm}
\scalebox{0.75}{%
\begin{tabular}{l rr rr rr}
\toprule
 & \multicolumn{2}{c}{Developed (DM)} & \multicolumn{2}{c}{Emerging (EM)} & \multicolumn{2}{c}{DM$-$EM} \\
\cmidrule(lr){2-3} \cmidrule(lr){4-5} \cmidrule(lr){6-7}
Surprise decile & $\mu$ & $\alpha$ & $\mu$ & $\alpha$ & $\Delta\mu$ & $\Delta\alpha$ \\
\midrule
\multicolumn{7}{c}{\textbf{Panel A:} Positive Surprises (Top Decile)} \\
\midrule
\texttt{sue\_q1} & 0.72 & 0.54 & 1.38 & 1.12 & $-$0.66 & $-$0.58 \\
 & (3.41) & (2.89) & (3.87) & (3.34) & ($-$2.14) & ($-$1.98) \\
\addlinespace
\texttt{sue\_q2} & 0.58 & 0.41 & 1.14 & 0.89 & $-$0.56 & $-$0.48 \\
 & (2.87) & (2.31) & (3.22) & (2.78) & ($-$1.89) & ($-$1.72) \\
\addlinespace
\texttt{car\_3d} & 0.63 & 0.48 & 1.21 & 0.97 & $-$0.58 & $-$0.49 \\
 & (3.12) & (2.67) & (3.45) & (3.01) & ($-$2.01) & ($-$1.84) \\
\addlinespace
\texttt{rev\_est} & 0.45 & 0.32 & 0.93 & 0.71 & $-$0.48 & $-$0.39 \\
 & (2.23) & (1.78) & (2.78) & (2.34) & ($-$1.67) & ($-$1.45) \\
\addlinespace
\midrule
\multicolumn{7}{c}{\textbf{Panel B:} Negative Surprises (Bottom Decile)} \\
\midrule
\texttt{sue\_q1} & $-$0.84 & $-$0.68 & $-$1.52 & $-$1.29 & 0.68 & 0.61 \\
 & ($-$3.67) & ($-$3.21) & ($-$4.12) & ($-$3.78) & (2.23) & (2.12) \\
\addlinespace
\texttt{sue\_q2} & $-$0.71 & $-$0.56 & $-$1.31 & $-$1.08 & 0.60 & 0.52 \\
 & ($-$3.12) & ($-$2.78) & ($-$3.56) & ($-$3.12) & (1.98) & (1.82) \\
\addlinespace
\texttt{car\_3d} & $-$0.78 & $-$0.63 & $-$1.45 & $-$1.22 & 0.67 & 0.59 \\
 & ($-$3.45) & ($-$3.01) & ($-$3.89) & ($-$3.45) & (2.18) & (2.01) \\
\addlinespace
\texttt{rev\_est} & $-$0.52 & $-$0.39 & $-$1.08 & $-$0.87 & 0.56 & 0.48 \\
 & ($-$2.34) & ($-$1.89) & ($-$3.01) & ($-$2.56) & (1.78) & (1.56) \\
\bottomrule
\end{tabular}
}
\end{center}
\end{table}

%%% --- TABLE 2: Summary Statistics / Descriptive (Archetype 12) --- %%%
%%% Clean descriptive statistics: no panels, no t-stats, right-aligned numerics.
%%% Pattern: l r r r r r r r, simple centering, footnotesize caption note.
\begin{table}[!ht]
\caption{[REMOVE] Summary Statistics.}
\begin{spacing}{1}
{\footnotesize
[REMOVE] This table reports descriptive statistics for the main variables. $N$ is the number of firm-quarter observations. P5, P95, and P99 denote the 5th, 95th, and 99th percentiles. Sample: 1990-Q1 to 2023-Q4.}
\end{spacing}
\vspace{-4mm}
\begin{center}
\label{tab:example_descriptive}
\vspace{2mm}
\scalebox{0.85}{%
\begin{tabular}{l rrrrrrr}
\toprule
Variable & $N$ & Mean & Median & SD & P5 & P95 & P99 \\
\midrule
SUE & 1{,}247{,}830 & 0.02 & 0.01 & 0.84 & $-$1.42 & 1.38 & 2.67 \\
\addlinespace
CAR [0,1] (\%) & 1{,}247{,}830 & 0.14 & 0.08 & 5.87 & $-$9.34 & 9.78 & 16.21 \\
\addlinespace
Mkt Cap (\$B) & 1{,}247{,}830 & 4.82 & 1.23 & 12.45 & 0.08 & 18.90 & 52.30 \\
\addlinespace
B/M & 1{,}183{,}450 & 0.71 & 0.54 & 0.62 & 0.09 & 1.78 & 3.12 \\
\addlinespace
Analyst coverage & 1{,}098{,}210 & 8.4 & 6.0 & 7.2 & 1.0 & 24.0 & 38.0 \\
\addlinespace
Bid-ask (\%) & 1{,}201{,}560 & 0.87 & 0.42 & 1.34 & 0.03 & 3.12 & 5.78 \\
\bottomrule
\end{tabular}
}
\end{center}
\end{table}

%%% --- TABLE 3: Lookup / Definition Table (Archetype 4/6) --- %%%
%%% Reference table with text columns: mnemonic, name, description, source.
%%% Pattern: l l p{6cm} l, no numeric data, no scaling needed.
\begin{table}[!ht]
\caption{[REMOVE] Variable Definitions.}
\vspace{-4mm}
\begin{center}
\label{tab:example_definitions}
\vspace{2mm}
\scalebox{0.80}{%
\begin{tabular}{l l p{6cm} l}
\toprule
Mnemonic & Name & Description & Source \\
\midrule
\texttt{sue} & Earnings surprise & Standardized unexpected earnings: actual EPS minus analyst consensus forecast, scaled by the stock price. & I/B/E/S \\
\addlinespace
\texttt{car\_3d} & Ann.\ return & Cumulative abnormal return over the three-day window [$-$1, +1] around the earnings announcement. & CRSP \\
\addlinespace
\texttt{drift\_60} & PEAD (60d) & Cumulative abnormal return over trading days [+2, +60] following the earnings announcement. & CRSP \\
\addlinespace
\texttt{mcap} & Market cap & Log of market capitalization at the end of the fiscal quarter (\$ millions). & Compustat \\
\addlinespace
\texttt{bm} & Book-to-market & Book value of equity divided by market capitalization, measured at fiscal year-end. & Compustat \\
\bottomrule
\end{tabular}
}
\end{center}
\end{table}

%% END:data


%%%%%%%%%%%%%%%%%%%%%%%%%%%%
% METHODOLOGY SECTION
%%%%%%%%%%%%%%%%%%%%%%%%%%%%

%% BEGIN:methodology
\section{Methodology}\label{sec:methodology}

[REMOVE] This section develops the econometric framework for measuring abnormal returns around earnings announcements. [REMOVE] We begin with a standard event-study specification and discuss the choice of benchmark model. [REMOVE] The abnormal return for firm $i$ on event day $\tau$ relative to its announcement date is:
\begin{equation}\label{eq:example_model}
  AR_{i,\tau} = r_{i,\tau} - \E[r_{i,\tau} \mid \mathcal{F}_{\tau-1}],
\end{equation}
where $r_{i,\tau}$ is the realized return and $\E[r_{i,\tau} \mid \mathcal{F}_{\tau-1}]$ is the expected return conditional on the pre-event information set. [REMOVE] We cumulate abnormal returns over the post-announcement window to obtain the drift measure, decomposed into an announcement component and a post-announcement component:
\begin{align}
  CAR_{i}[0,1]  &= \sum_{\tau=0}^{1} AR_{i,\tau}, \label{eq:example_return} \\
  PEAD_{i}[2,T] &= \sum_{\tau=2}^{T} AR_{i,\tau}, \label{eq:example_signal}
\end{align}
where $CAR_{i}[0,1]$ captures the immediate price reaction and $PEAD_{i}[2,T]$ captures the subsequent drift over $T$ trading days. [REMOVE] The central question is whether $\E[PEAD_{i}[2,T] \mid SUE_{i}] \neq 0$, which would indicate that prices do not fully incorporate the information in earnings surprises at the announcement date (see Eq.~\eqref{eq:example_model}).

%% END:methodology


%%%%%%%%%%%%%%%%%%%%%%%%%%%%
% RESULTS SECTION
%%%%%%%%%%%%%%%%%%%%%%%%%%%%

%% BEGIN:results
\section{Results}\label{sec:results}

[REMOVE] This section presents the main empirical findings. [REMOVE] We first document substantial cross-country variation in the PEAD and then decompose it into long-leg and short-leg contributions. [REMOVE] Table~\ref{tab:example_comparison} reports drift magnitudes under alternative benchmark models, and Figure~\ref{fig:example} displays cumulative abnormal return paths for the top and bottom surprise deciles.

%%% --- TABLE 4: Multi-Leg Decomposition (Archetype 2) --- %%%
%%% Long/Short/Long-Short breakdown: wide table with 9+ columns.
%%% Pattern: \scalebox{0.72}, 3 leg groups with \multicolumn + \cmidrule, t-stats on sub-rows.
\begin{table}[!ht]
\caption{[REMOVE] PEAD Decomposition by Portfolio Leg and Announcement Window (\%).}
\begin{spacing}{1}
{\footnotesize
[REMOVE] This table decomposes the PEAD strategy into long-leg (top SUE decile) and short-leg (bottom decile) contributions across three windows: [0,1] announcement, [2,20] short-horizon drift, and [21,60] medium-horizon drift. $t$-statistics (Newey-West) in parentheses. Sample: 1990-Q1 to 2023-Q4.}
\end{spacing}
\vspace{-4mm}
\begin{center}
\label{tab:example_decomposition}
\vspace{2mm}
\scalebox{0.72}{%
\begin{tabular}{l rrr rrr rrr}
\toprule
 & \multicolumn{3}{c}{Long (Top SUE, \%)} & \multicolumn{3}{c}{Short (Bot SUE, \%)} & \multicolumn{3}{c}{Long$-$Short (\%)} \\
\cmidrule(lr){2-4} \cmidrule(lr){5-7} \cmidrule(lr){8-10}
Window & Gross & Net & Costs & Gross & Net & Costs & Gross & Net & Costs \\
\midrule
{[0,1]} & 1.82 & 1.71 & 0.11 & $-$1.94 & $-$1.83 & $-$0.11 & 3.76 & 3.54 & 0.22 \\
 & (8.34) & (7.89) & (6.12) & ($-$8.78) & ($-$8.34) & ($-$6.45) & (12.01) & (11.34) & (7.89) \\[1mm]
{[2,20]} & 0.68 & 0.41 & 0.27 & $-$0.74 & $-$0.47 & $-$0.27 & 1.42 & 0.88 & 0.54 \\
 & (3.45) & (2.12) & (4.56) & ($-$3.67) & ($-$2.34) & ($-$4.78) & (5.01) & (3.12) & (6.34) \\[1mm]
{[21,60]} & 0.34 & 0.12 & 0.22 & $-$0.41 & $-$0.19 & $-$0.22 & 0.75 & 0.31 & 0.44 \\
 & (1.89) & (0.67) & (3.78) & ($-$2.12) & ($-$0.98) & ($-$3.89) & (2.78) & (1.15) & (5.12) \\
\bottomrule
\end{tabular}
}
\end{center}
\end{table}

%%% --- TABLE 5: Compact Comparison (Archetype 3) --- %%%
%%% Mean returns and alphas side-by-side with 2 panels.
%%% Pattern: \scalebox{0.78}, l ccc ccc, \multicolumn panel labels flush-left.
\begin{table}[!ht]
\caption{[REMOVE] PEAD Under Alternative Benchmark Models (Excess Returns).}
\begin{spacing}{1}
{\footnotesize
[REMOVE] This table reports the PEAD over [+2, +60] days under three benchmark models: market model (MM), Fama-French three-factor (FF3), and Fama-French five-factor (FF5). $t$-statistics (Newey-West, lags $= \lfloor T^{0.25} \rfloor$) in parentheses. Sample: 1990-Q1 to 2023-Q4.}
\end{spacing}
\vspace{-4mm}
\begin{center}
\label{tab:example_comparison}
\vspace{2mm}
\scalebox{0.78}{%
\begin{tabular}{l ccc ccc}
\toprule
 & \multicolumn{3}{c}{Top Decile (\%)}  & \multicolumn{3}{c}{Bottom Decile (\%)} \\
\cmidrule(lr){2-4} \cmidrule(lr){5-7}
 & MM & FF3 & FF5 & MM & FF3 & FF5 \\
\midrule
\multicolumn{7}{l}{\textbf{Panel A:} Developed Markets} \\
\midrule
\texttt{sue\_q1} & 1.14 & 0.92 & 0.78 & $-$1.28 & $-$1.04 & $-$0.89 \\
 & (4.56) & (3.89) & (3.34) & ($-$5.01) & ($-$4.34) & ($-$3.78) \\
\texttt{sue\_q2} & 0.98 & 0.78 & 0.65 & $-$1.08 & $-$0.86 & $-$0.72 \\
 & (3.89) & (3.21) & (2.78) & ($-$4.23) & ($-$3.56) & ($-$3.01) \\
\texttt{car\_3d} & 1.06 & 0.84 & 0.71 & $-$1.18 & $-$0.95 & $-$0.81 \\
 & (4.12) & (3.45) & (2.98) & ($-$4.56) & ($-$3.89) & ($-$3.34) \\
\midrule
\multicolumn{7}{l}{\textbf{Panel B:} Emerging Markets} \\
\midrule
\texttt{sue\_q1} & 2.34 & 1.98 & 1.72 & $-$2.56 & $-$2.18 & $-$1.89 \\
 & (5.67) & (4.89) & (4.23) & ($-$6.12) & ($-$5.34) & ($-$4.67) \\
\texttt{sue\_q2} & 1.89 & 1.56 & 1.34 & $-$2.12 & $-$1.78 & $-$1.52 \\
 & (4.56) & (3.89) & (3.34) & ($-$5.01) & ($-$4.23) & ($-$3.67) \\
\bottomrule
\end{tabular}
}
\end{center}
\end{table}

%%% --- FIGURE 1: Example Figure --- %%%
%%% Standard figure with caption below, footnotesize notes in \justify + \spacing{1}.
%%% Pattern: \begin{figure}[h!], placeholder graphic, negative \vspace, \label at end.
\begin{figure}[h!]
\begin{center}
\fbox{\rule{0pt}{4cm}\rule{0.95\textwidth}{0pt}}
\end{center}
\vspace{-4mm}
\caption{[REMOVE] Cumulative abnormal returns by earnings surprise decile.}
\vspace{-2mm}
\begin{justify}
\begin{spacing}{1}
\footnotesize{
[REMOVE] Shows cumulative abnormal returns (CARs) over trading days [$-$5, +60] relative to the earnings announcement date. Panel~(A) plots the top and bottom SUE deciles for developed markets. Panel~(B) plots the same for emerging markets. The $y$-axis is in percentage points. Returns are equal-weighted. Sample: 1990-Q1 to 2023-Q4.
}
\end{spacing}
\end{justify}
\vspace{-12mm}
\label{fig:example}
\end{figure}

%% END:results


%%%%%%%%%%%%%%%%%%%%%%%%%%%%
% CONCLUSION SECTION
%%%%%%%%%%%%%%%%%%%%%%%%%%%%

%% BEGIN:conclusion
\section{Conclusion}\label{sec:conclusion}

[REMOVE] This paper documents the post-earnings announcement drift across 40 international equity markets and quantifies its economic significance net of transaction costs. [REMOVE] The drift is larger and more persistent in emerging markets, where analyst coverage is sparser and short-selling constraints are tighter. [REMOVE] In developed markets, the net-of-cost drift has declined over time, consistent with increased arbitrage activity following the documentation of the anomaly \citep{FamaFrench_2015}. [REMOVE] These findings highlight the importance of accounting for market microstructure when evaluating the profitability of anomaly-based strategies.

%% END:conclusion


%%%%%%%%%%%%%%%%%%%%%%%%%%%%
% BIBLIOGRAPHY
%%%%%%%%%%%%%%%%%%%%%%%%%%%%
\clearpage
\bibliographystyle{chicago}
\bibliography{template_references} % %%BIB_FILE%% -- skill replaces this line


%%%%%%%%%%%%%%%%%%%%%%%%%%%%
% APPENDIX
%%%%%%%%%%%%%%%%%%%%%%%%%%%%
\clearpage
\newpage
\appendix

\begin{center}
 {\LARGE \textbf{Appendix}}
\end{center}

% Change section numbering style
\renewcommand{\theequation}{A.\arabic{equation}}%
\renewcommand{\thefigure}{A.\arabic{figure}} \setcounter{figure}{0}
\renewcommand{\thetable}{A.\arabic{table}} \setcounter{table}{0}


%% BEGIN:appendix-a
%%%%%%%%%%%%%%%%%%%%%%%%%%%%
% APPENDIX A
%%%%%%%%%%%%%%%%%%%%%%%%%%%%
\vspace{-.25cm}
\section{Additional Results} \label{app:a}

[REMOVE] This appendix provides supplementary analyses. [REMOVE] Table~\ref{tab:example_summary} summarizes results across all signal categories.

%%% --- TABLE 6: Appendix Summary Table (Archetype 11) --- %%%
%%% Cross-factor summary with counts and aggregate statistics.
%%% Pattern: Simple centering, no scalebox, l c c c c columns.
\begin{table}[!ht]
\caption{[REMOVE] PEAD by Country Group and Surprise Measure.}
\begin{spacing}{1}
{\footnotesize
[REMOVE] This table summarizes the 60-day PEAD across country groups and earnings surprise measures. Count is the number of country-quarter observations. \% Sig.\ is the fraction with $|t| > 1.96$. $\overline{|CAR|}$ is the mean absolute cumulative abnormal return (\%). $\overline{|t|}$ is the mean absolute $t$-statistic.}
\end{spacing}
\vspace{-4mm}
\begin{center}
\label{tab:example_summary}
\vspace{2mm}
\begin{tabular}{l cccc}
\toprule
Country group & Count & \% Sig. & $\overline{|CAR|}$ & $\overline{|t|}$ \\
\midrule
North America & 14 & 71\% & 1.42 & 3.87 \\
\addlinespace
Europe & 18 & 56\% & 1.18 & 2.98 \\
\addlinespace
Asia-Pacific & 12 & 67\% & 1.65 & 3.45 \\
\addlinespace
Emerging & 16 & 81\% & 2.34 & 4.12 \\
\midrule
All & 60 & 68\% & 1.62 & 3.56 \\
\bottomrule
\end{tabular}
\end{center}
\end{table}

%%% [REMOVE] Longtable example — delete when writing your paper
% \begin{longtable}{l p{10cm}}
% \caption{\textbf{Variable definitions.}}\label{tab:vardefs}\\
% \toprule Variable & Definition \\ \midrule
% \endfirsthead
% \multicolumn{2}{c}{\footnotesize\emph{Table~\thetable{} continued}} \\
% \toprule Variable & Definition \\ \midrule
% \endhead
% \midrule \multicolumn{2}{r}{\footnotesize\emph{Continued on next page}} \\
% \endfoot
% \bottomrule
% \endlastfoot
% HML & High minus low book-to-market portfolio return \\
% SMB & Small minus big size portfolio return \\
% \end{longtable}
% {\footnotesize Notes go here --- OUTSIDE longtable. Do not use
% \texttt{\textbackslash begin\{spacing\}} inside longtable.}

%% END:appendix-a


%%%%%%%%%%%%%%%%%%%%%%%%%%%%
%%% Internet Appendix %%%%
%%%%%%%%%%%%%%%%%%%%%%%%%%%%

%% BEGIN:internet-appendix

\renewcommand{\theequation}{IA.\arabic{equation}}%
\renewcommand{\thefigure}{IA.\arabic{figure}} \setcounter{figure}{0}
\renewcommand{\thetable}{IA.\Roman{table}} \setcounter{table}{0}
\renewcommand{\thesection}{IA.\arabic{section}} \setcounter{section}{0}
\setcounter{footnote}{0}

\pagebreak
\clearpage
\setcounter{page}{1}

\begin{center}
{\large \textbf{Internet Appendix for:}}

\medskip

{\LARGE \textbf{[REMOVE] Your Paper Title Here}}

\bigskip

\normalsize{\textbf{Abstract}}
\end{center}

\noindent [REMOVE] This Internet Appendix provides additional information, tables, figures, and empirical results supporting the main text.

\newpage

\singlespacing

%%%%%%%%%%%%%%%%%%%%%%%%%%%%
% IA SECTION 1
%%%%%%%%%%%%%%%%%%%%%%%%%%%%
\section{[REMOVE] PEAD by Firm Characteristics}\label{sec:IA-desc-stats}

[REMOVE] This section reports the PEAD conditional on firm-level characteristics. [REMOVE] Table~\ref{tab:ia_example} presents results across size and analyst coverage terciles.

%%% --- IA TABLE: Results with Roman numeral numbering --- %%%
%%% Internet Appendix tables use Roman numerals (IA.I, IA.II, etc.)
%%% and IA-prefixed section numbers (IA.1, IA.2, etc.)
\begin{table}[!ht]
\caption{[REMOVE] PEAD [+2, +60] by Size and Analyst Coverage Tercile (\%).}
\begin{spacing}{1}
{\footnotesize
[REMOVE] This table reports the 60-day PEAD across size and analyst coverage terciles. Small, medium, and large refer to terciles of market capitalization. $t$-statistics (Newey-West, lags $= \lfloor T^{0.25} \rfloor$) in parentheses.}
\end{spacing}
\vspace{-4mm}
\begin{center}
\label{tab:ia_example}
\vspace{2mm}
\scalebox{0.80}{%
\begin{tabular}{l rr rr rr}
\toprule
 & \multicolumn{2}{c}{Small} & \multicolumn{2}{c}{Medium} & \multicolumn{2}{c}{Large} \\
\cmidrule(lr){2-3} \cmidrule(lr){4-5} \cmidrule(lr){6-7}
Coverage & Top SUE & Bot SUE & Top SUE & Bot SUE & Top SUE & Bot SUE \\
\midrule
Low & 2.14 & $-$2.38 & 1.45 & $-$1.62 & 0.89 & $-$1.01 \\
 & (4.56) & ($-$4.89) & (3.21) & ($-$3.45) & (2.12) & ($-$2.34) \\
\addlinespace
Medium & 1.67 & $-$1.89 & 1.12 & $-$1.28 & 0.67 & $-$0.78 \\
 & (3.78) & ($-$4.12) & (2.78) & ($-$3.01) & (1.78) & ($-$1.98) \\
\addlinespace
High & 1.23 & $-$1.41 & 0.78 & $-$0.92 & 0.42 & $-$0.51 \\
 & (2.89) & ($-$3.23) & (2.01) & ($-$2.34) & (1.23) & ($-$1.45) \\
\bottomrule
\end{tabular}
}
\end{center}
\end{table}

%%%%%%%%%%%%%%%%%%%%%%%%%%%%
% IA SECTION 2
%%%%%%%%%%%%%%%%%%%%%%%%%%%%
\section{[REMOVE] Alternative Earnings Surprise Measures}\label{sec:IA-robustness}

[REMOVE] This section evaluates the robustness of the PEAD to alternative definitions of the earnings surprise. [REMOVE] We consider analyst forecast errors from I/B/E/S, seasonal random walk forecasts, and management guidance-based surprises; results are qualitatively similar across all three measures.

%% END:internet-appendix

\end{document}
